\section{Logs and Output Analysis}
\label{sec:analysis}

\subsection{The anomalous EN→DA zero-shot prediction file}

The most significant pipeline anomaly is the English-to-Danish BERT zero-shot run (\texttt{english\_to\_danish\_zero\_shot}). Its development prediction file contains 2{,}001 sentences, while the Danish development set has only 564. The strict F1 against the Danish gold file is 0.3\% (effectively zero), yet the saved training-time metric is 79.3 F1. The most probable cause is a stale or incorrect path in the evaluation call during training: the trainer evaluated the model on the English development set rather than the Danish one, saving the English-set performance as the Danish dev metric. The model was trained on English data, so evaluating it on English would produce a plausible score; evaluating the same predictions against Danish gold produces near-zero because the sentence count mismatches and the underlying predictions are English-language.

\textbf{Blast-radius check.} To assess whether the same bug could silently affect other runs, we compared per-experiment sentence counts between prediction files and the correct gold splits. Only \texttt{english\_to\_danish\_zero\_shot} has a mismatch on the development set (2{,}001 vs 564). All other 16 experiments produce prediction files with the correct sentence count. We also inspected the metric gap (strict F1 vs saved eval F1) for every other run: apart from this artifact, all gaps are below 2 F1 points, which is within expected scoring-method variation. We conclude that the bug affected only this one run and did not propagate to the fine-tuned model that starts from its checkpoint. The checkpoint itself is a valid English NER model; the fine-tuning step provides its own ground truth by continuing training on Danish data.

On the test split, the same experiment's prediction file does \emph{not} have a sentence count mismatch, producing a test strict F1 of 78.8. This is consistent with the stored metric, suggesting the test evaluation used the correct file. We exclude this run from transfer comparisons but report the test score for completeness.

\subsection{MISC entities are absent from evaluation}

Per-entity-type results show 0.0 F1 for MISC across all systems on all splits. This is not a model failure. The Danish Universal NER file uses a two-column annotation format: the primary label column (column 3, index 2) contains the DDT-style annotations (PER, LOC, ORG, O), while a secondary column encodes a separate annotation layer that includes MISC labels (\texttt{B-MISC\#O}). Both training code and the evaluation script (\texttt{span\_f1.py}) read from the primary column, in which MISC does not appear. Models therefore never see MISC during training and there are no MISC entities to predict in the evaluation column. The 32 MISC annotations in the Danish dev file and 44 in the test file reside in the secondary column only.

\subsection{Metric consistency for well-behaved runs}

Excluding the artifact, the remaining 16 runs show tight agreement between strict span F1 (computed by \texttt{span\_f1.py} on the saved prediction files) and the evaluation F1 logged during training (Figure~\ref{fig:strict_eval_scatter}). The largest discrepancy is 1.2 F1 points (Danish XLM-R 25\%), well within the range expected from different tokenisation-boundary handling between the two evaluation paths. This cross-check confirms that the training pipeline reported consistent and reproducible scores throughout.
